% Graphical abstract: what the paper does.
\begin{figure*}[!t]
\centering
\resizebox{\textwidth}{!}{%
\begin{tikzpicture}[
  font=\small,
  hdr/.style={font=\small\bfseries, align=center},
  ax/.style={draw=black!60, rounded corners=2pt, fill=blue!7,
             text width=34mm, align=center, minimum height=8mm, inner sep=3pt},
  lay/.style={draw=black!60, rounded corners=2pt, fill=teal!9,
              text width=34mm, align=center, minimum height=8mm, inner sep=3pt},
  fail/.style={draw=orange!70!black, rounded corners=2pt, fill=orange!12,
               text width=34mm, align=center, minimum height=8mm, inner sep=3pt},
  big/.style={-{Stealth[length=4mm]}, line width=1.6pt, draw=black!40}
]

% ---- column 1: five classification axes
\node[hdr] (h1) at (0,1.35) {How systems are classified};
\node[ax] (a1) at (0,0.35)   {Direction of knowledge flow};
\node[ax] (a2) at (0,-0.85)  {Lifecycle stage};
\node[ax] (a3) at (0,-2.05)  {Coupling depth};
\node[ax] (a4) at (0,-3.25)  {Knowledge object};
\node[ax] (a5) at (0,-4.45)  {Evidence maturity};

% ---- column 2: seven evaluation layers
\node[hdr] (h2) at (6.2,1.35) {Where benefit is measured};
\node[lay] (l1) at (6.2,0.55)   {Graph and schema};
\node[lay] (l2) at (6.2,-0.45)  {Linking and query};
\node[lay] (l3) at (6.2,-1.45)  {Retrieval};
\node[lay] (l4) at (6.2,-2.45)  {Evidence use};
\node[lay] (l5) at (6.2,-3.45)  {Generation};
\node[lay] (l6) at (6.2,-4.45)  {Systems cost};
\node[lay] (l7) at (6.2,-5.45)  {Human utility};

% ---- column 3: failure modes
\node[hdr] (h3) at (12.4,1.35) {Why it fails in practice};
\node[fail] (f1) at (12.4,0.35)  {Coverage gaps, stale edges};
\node[fail] (f2) at (12.4,-0.85) {Entity-linking error};
\node[fail] (f3) at (12.4,-2.05) {Context dilution};
\node[fail] (f4) at (12.4,-3.25) {Unsupported citation};
\node[fail] (f5) at (12.4,-4.45) {Circular self-validation};

% ---- flow arrows between columns
\draw[big] (2.05,-2.05) -- (4.15,-2.05);
\draw[big] (8.25,-2.05) -- (10.35,-2.05);

% ---- conclusion strip
\node[draw=black!55, fill=blue!5, rounded corners=2pt, text width=168mm,
      align=center, inner sep=5pt, font=\small] at (6.2,-7.1)
  {\textbf{Bounded finding.} A graph helps when it holds knowledge the model lacks, entity linking
   succeeds, retrieval returns a compact and relevant subgraph, and the comparison is a matched
   text-retrieval baseline. Otherwise it adds cost, and sometimes error.};
\end{tikzpicture}%
}
\caption{The structure of this survey's argument. \textbf{Left:} systems are classified along five axes rather than by direction of knowledge flow alone, because pasting triples into a prompt and training a graph-aware encoder are different interventions that the usual taxonomy treats alike. Coupling depth runs from C0, where graph and model share an application but exchange nothing structured, to C5, where the graph constrains decoding and may receive controlled write-back; evidence maturity runs from E0, a prototype with no systematic evaluation, to E6, governed routine use. \textbf{Centre:} benefit is measured at seven separable layers, L1 through L7 in the order shown, so that a failure can be attributed to the graph, the linker, the retriever, or the generator rather than to the system as a whole. \textbf{Right:} the failure modes that recur across the reviewed systems. Circular self-validation arises only in closed-loop designs that write model output back into the graph, and it has no analogue in unidirectional systems.}
\label{fig:abstract}
\end{figure*}
